\documentclass[11pt]{amsart}

\usepackage[T1]{fontenc}
\usepackage{lmodern}
\usepackage{microtype}
\usepackage{amsmath,amssymb,amsthm}
\usepackage[hidelinks]{hyperref}

\hypersetup{
  pdftitle={Edge-Deletion Counterexamples to the Aliabadi--Krop Sharp Uniform-Gap Formula},
  pdfsubject={Perfect matchings and the Alon--Friedland bound},
  pdfkeywords={perfect matching, Alon--Friedland bound, extremal graph, uniform gap}
}

\newtheorem{theorem}{Theorem}
\newtheorem{corollary}[theorem]{Corollary}
\theoremstyle{remark}
\newtheorem{remark}[theorem]{Remark}

\newcommand{\perfmat}{\operatorname{perfmat}}

\title[Edge-deletion counterexamples]
{Edge-Deletion Counterexamples to the Aliabadi--Krop
Sharp Uniform-Gap Formula}

\date{}
\subjclass[2020]{05C70, 05C35}
\keywords{perfect matching, Alon--Friedland bound, extremal graph, uniform gap}

\begin{document}

\begin{abstract}
Aliabadi and Krop asked for the sharp normalized perfect-matching count
below the Alon--Friedland bound among non-extremal graphs of maximum degree
at most $\Delta$. For every $\Delta\ge4$, we show that deleting one edge
from $K_{\Delta,\Delta}$ gives a ratio larger than both candidates in their
proposed formula. The same family yields an upper bound of order
$(\log\Delta)/\Delta^2$ on any uniform gap constant, improving their
order-$1/\Delta$ obstruction.
\end{abstract}

\maketitle

\section{The counterexamples}

For a finite simple graph $G$ of even order and without isolated vertices, set
\[
U(G)=\prod_{v\in V(G)}\bigl(d_G(v)!\bigr)^{1/(2d_G(v))}.
\]
The Alon--Friedland bound states that
$\perfmat(G)\le U(G)$, with equality exactly for disjoint unions of
balanced complete bipartite graphs \cite{AlonFriedland}.

For $\Delta\ge3$, Aliabadi and Krop define
\[
R_\Delta=\sup\left\{
\frac{\perfmat(G)}{U(G)}:
\begin{array}{l}
G\text{ has even order, no isolated vertices, and }\Delta(G)\le\Delta,\\
G\text{ is not a disjoint union of balanced complete bipartite graphs}
\end{array}
\right\}
\]
and ask in \cite[Problem~4.9]{AliabadiKrop} whether
\begin{equation}\label{eq:AK}
R_\Delta=
\max\left\{\left(\frac34\right)^{1/3},\rho_{\Delta-1}\right\},
\qquad
\rho_r=\frac{(r!)^{1/(r(r+1))}}{(r+1)^{1/(r+1)}}.
\end{equation}

For $r\ge2$, let $H_r=K_{r,r}-e$, where $e$ is one edge, and write
\[
\sigma_r=\frac{\perfmat(H_r)}{U(H_r)}.
\]
Aliabadi and Krop observed $H_3$ as the runner-up in their enumeration of
connected graphs on at most seven vertices \cite[Remark~4.10]{AliabadiKrop}.
The full deletion family refutes \eqref{eq:AK} from $r=4$ onward.

\begin{theorem}\label{thm:main}
For every integer $r\ge4$, the graph $H_r$ is admissible in the definition
of $R_r$ and
\[
\sigma_r=
\frac{r-1}
{r^{(r-1)/r}\bigl((r-1)!\bigr)^{1/(r(r-1))}}
>
\max\left\{\left(\frac34\right)^{1/3},\rho_{r-1}\right\}.
\]
Consequently, the formula \eqref{eq:AK} fails for every $\Delta\ge4$.
\end{theorem}

\begin{proof}
The graph $H_r$ is connected, has even order, has no isolated vertices, and has
maximum degree $r$. Since one cross-edge is missing, it is not a balanced
complete bipartite graph. Of the $r!$ perfect matchings of $K_{r,r}$,
exactly $(r-1)!$ use the deleted edge. Hence
\[
\perfmat(H_r)=(r-1)(r-1)!.
\]
There are $2r-2$ vertices of degree $r$ and two of degree $r-1$, so
\[
U(H_r)=(r!)^{(r-1)/r}\bigl((r-1)!\bigr)^{1/(r-1)},
\]
which gives the stated formula for $\sigma_r$.

By strict arithmetic--geometric mean,
\begin{equation}\label{eq:amgm}
\bigl((r-1)!\bigr)^{1/(r-1)}<\frac r2.
\end{equation}
Therefore
\[
\frac{\sigma_r}{\rho_{r-1}}
=
\frac{r-1}
{r^{(r-2)/r}\bigl((r-1)!\bigr)^{2/(r(r-1))}}
>
\left(1-\frac1r\right)2^{2/r}>1.
\]
For the last inequality, the function
$x\mapsto x\log(1-1/x)$ is increasing on $(1,\infty)$, since its derivative
is
\[
\frac1{x-1}-\log\left(1+\frac1{x-1}\right)>0.
\]
Thus, for $r\ge3$,
\[
\left[\left(1-\frac1r\right)2^{2/r}\right]^r
=4\left(1-\frac1r\right)^r
\ge4\left(\frac23\right)^3>1.
\]

It remains to compare $\sigma_r$ with $(3/4)^{1/3}$. From
\eqref{eq:amgm},
\[
\sigma_r>
q(r):=\left(1-\frac1r\right)2^{1/r}.
\]
The function $q$ is increasing on $(1,\infty)$ because
\[
\frac{d}{dx}\log q(x)
=
\frac{x/(x-1)-\log2}{x^2}>0.
\]
Moreover,
\[
\left(\frac{q(5)}{(3/4)^{1/3}}\right)^{15}
=
\frac{2^{43}}{3^5 5^{15}}>1,
\]
where the final inequality follows from $3^5<2^8$ and $5^3<2^7$.
Hence $\sigma_r>(3/4)^{1/3}$ for every $r\ge5$. For $r=4$, direct
simplification gives
\[
\left(\frac{\sigma_4}{(3/4)^{1/3}}\right)^{12}
=
\frac{3^7}{2^{11}}>1.
\]
This proves the two strict comparisons. Taking $r=\Delta$ completes the
proof.
\end{proof}

\section{A sharper asymptotic obstruction}

Aliabadi and Krop obtained
\[
c(\Delta)\le \frac1\Delta+O\left(\frac{\log\Delta}{\Delta^2}\right)
\]
from an added-edge construction \cite[Corollary~4.8]{AliabadiKrop}.
The deletion family gives the following stronger obstruction.

\begin{corollary}\label{cor:gap}
As $\Delta\to\infty$,
\[
1-R_\Delta
\le 1-\sigma_\Delta
=
\frac{\log(2\pi\Delta)-1}{2\Delta^2}
+O\left(\frac{\log\Delta}{\Delta^3}\right).
\]
In particular, any constant $c(\Delta)$ valid in
\cite[Problem~4.4]{AliabadiKrop} satisfies the same upper bound.
\end{corollary}

\begin{proof}
From the formula in Theorem~\ref{thm:main},
\[
\log\sigma_r
=
\log\left(1-\frac1r\right)
+\frac{\log r}{r}
-\frac{\log\Gamma(r)}{r(r-1)}.
\]
Stirling's formula gives
\[
\log\sigma_r
=-\frac{\log(2\pi r)-1}{2r^2}
+O\left(\frac{\log r}{r^3}\right).
\]
The square of the main term is
$o((\log r)/r^3)$, so exponentiation yields
\[
1-\sigma_r
=
\frac{\log(2\pi r)-1}{2r^2}
+O\left(\frac{\log r}{r^3}\right).
\]
Since $R_\Delta\ge\sigma_\Delta$, the result follows.
\end{proof}

\begin{remark}
The theorem does not determine $R_\Delta$, settle the case $\Delta=3$, or
resolve whether $R_\Delta<1$ for each fixed $\Delta$.
\end{remark}

\begin{thebibliography}{9}

\bibitem{AlonFriedland}
N.~Alon and S.~Friedland,
\emph{The maximum number of perfect matchings in graphs with a given degree sequence},
Electron. J. Combin. \textbf{15} (2008), no.~1, Note~N13, 2~pp.,
\href{https://doi.org/10.37236/888}{doi:10.37236/888}.

\bibitem{AliabadiKrop}
M.~Aliabadi and E.~Krop,
\emph{Odd-cycle defects in the Alon--Friedland bound},
\href{https://arxiv.org/abs/2607.14134v1}
{arXiv:2607.14134v1 [math.CO]}, 2026.

\end{thebibliography}

\end{document}